\documentclass[english,twoside, openright, hidelinks, 11 pt, class=report,crop=false]{standalone}
\input{../../preamb}



\begin{document}
{\Large Eksamen R1 Høsten 2025\hfill {\footnotesize Løsning fra \color{blue} \href{https://openmath.online}{OpenMathBooks}}}	

\subsection*{Oppgave 1}
\abc{
\item $f(x)=3x^2+\frac{1}{2}x^{-\frac{1}{2}}$
\item Vi setter $u = 2x-3$ og $v=e^x$. Da er $u'=2$ og $v'=e^x$. Av divisjonsformelen ved derivasjon har vi at
\[ f'(x)=\left(\frac{u}{v}\right)'=\frac{u'v-uv'}{v^2}=\frac{2e^x-(2x-3)e^x}{\left(e^x\right)^2}=\frac{5-2x}{e^x} \]
\item $g'(2)=\frac{5-2\cdot2}{e^2}=\frac{1}{e^2}$ og $g'(3)=\frac{5-2\cdot3}{e^3}=-\frac{1}{e^3}$
\item Grafen går fra å være stigende til synkende. Dette betyr også at $g$ har en lokalt toppunkt på intervallet.
}

\subsection*{Oppgave 2}
\abc{
\item Ved å sette $u=\lg x$, kan vi skrive
\[ u^2-2u-8=0 \]
Av \textit{abc}-formelen har vi at
\begin{align*}
	u &= \frac{-(-2)\pm\sqrt{(-2)^2-4\cdot1\cdot(-8)}}{2\cdot1} \br
	&=\frac{2\pm\sqrt{36}}{2} \br
	&= \frac{2\pm 6}{2}
\end{align*}
Dermed er
\[ u=4\qquad\vee\qquad u=-2 \]
Antatt at $\lg$ har 10 som base er
\[ x = 10^{4}\qquad\vee\qquad x=10^{-2} \]
\item Vi har at $\sqrt[3]{64}=\sqrt[3]{4^3}=4$. Dermed er $\frac{1}{64}=4^{-3}$, og da er 
\[ \log_4 \frac{1}{64}=-3 \]
}

\subsection*{Oppgave 3}
\abc{
\item Nevneren til uttrykket går mot 0 når $x$ går mot $-2$, mens telleren går mot 14. Dermed eksisterer ikke grenseverdien.
\item 1) Grensen eksisterer hvis $x+2$ er en faktor i telleren. For en $a=2+b$, skal vi da kunne skrive telleren som $x^2+(2+b)x+2b$. Dette må bety at $b=1$,og da er $a=3$. \os
2) Da $2\cdot(-4)=-8$ og $2+(-4)=-2$, har vi at
\[ \lim\limits_{x\to-2} \frac{x^2+3x+2}{x^2-2x-8}=\lim\limits_{x\to-2} \frac{(x+2)(x+1)}{(x+2)(x-4)}=\frac{-2+1}{-2-4}=\frac{1}{6} \] 
}

\subsection*{Oppgave 4}
\abc{
\item $\vv{AB}=[3-(-2), 2-3]=[5, -1]$. Videre er
\[ \left|\vv{AB}\right|=\sqrt{5^2+(-1)^2}=\sqrt{26} \]
\item Parameteriseringen av linja kan vi skrive som
\[ \vec{r}(t)=(-2, 3)+t[5, -1] \]
Linja skjærer $x$-aksen når $y$-verdien til $\vec{r}$ er lik 0, som er når $t=3$. Dermed er
\[ C=\vec{r}(3)=(-2+15, 0)=(-13, 0) \]
\item $\vv{BD}=[-1, t-2]$. Hvis $\angle ABD=90^\circ$, har vi at
\begin{align*}
	\vv{AB}\cdot\vv{BD}&=0 \\
	[5, -1]\cdot[-1, t-2]&=0 \\
	5\cdot(-1)+(-1)(t-2)&=0 \\
	-5-t+2&=0 \\
	t&=-3
\end{align*}
Dermed er $D=(2, -3)$.
}

\subsection*{Oppgave 5}
\abc{
\item Vi setter $u=4x^2$ og $\ln x=v$. Da er $u'=8x$ og $v'=\frac{1}{x}$. Av derivasjonsregelen ved multiplikasjon har vi at
\[ f'(x)=u'v+uv'=8x\ln x+4x^2\cdot\frac{1}{x}=4x(2\ln x+1) \]
$f$ har toppunkt/bunnpunkt der hvor $f'(x)=0$. $4x=0$ når $x=0$, men $\ln 0$ er ikke definert. Dette betyr at $f'(x)=0$ bare når $2\ln x=-1$, som betyr at $x=e^{-\frac{1}{2}}$. $4x$ er positiv på intervallet for alle $x>0$, mens $2\ln x +1$ skifter fortegn fra negativ til positiv når $x=e^{-\frac{1}{2}}$. Dette er altså et minimumspunkt for $f$. Vi har at
\[ f\left(e^{-\frac{1}{2}}\right)=4\cdot \left(e^{-\frac{1}{2}}\right)^2\cdot \ln e^{-\frac{1}{2}}=4e^{-1}\cdot\left(-\frac{1}{2}\right)=-2e^{-1} \]
Altså er $\left(e^{-\frac{1}{2}}, -2e^{-1}\right)$ et bunnpunkt på grafen til $f$.
\item Eleven prøver å finne nullpunktet til funksjonen. \os

Vi lar $x=x_1$ være løsningen av likningen $f(x)=0$.
\\Linje 11: Starter med "å gjette" at $x_1$ er midtpunktet mellom $a$ og $b$.\\
Linje 13: Itererer så lenge $f(m)$ er over $0,0001$.\\
Linje 15-18: Setter $b=m$ hvis $f(a)\cdot f(m)<0$, og $a=m$ hvis ikke. Hvis vilkåret er oppfylt, har $f(a)$ og $f(m)$ forskjellig fortegn, og siden $f(a)<0$, vet vi at $x_1\in[a, m]$. Da må ny høyre grense settes til $m$. Hvis vilkåret  ikke er oppfylt er $x_1\in[m, b]$. Da må ny venstre grense settes til $m$.\\
Linje 20: ''Gjetter'' at $x$ er midpunktet mellom nye $a$ og $b$. \os

$f(x)=0$ når $\ln x =0$, altså når $x=1$.
}

\newpage
\subsection*{Oppgave 1}
\abc{
\item Lager en liste med punkt, og bruker \texttt{RegLogist}. Får da funksjonen $f(x)$ som samsvarer med $F(t)$ i oppgaven. Vi ser at grafen ligger rimelig nærme punktene, som tyder på at modellen er gyldig på intervallet $[0, 20]$. Da er det rimelig å tro at modellen vil være gyldig også for noen år før 1960 og for noen år etter 1980. Men når $t$ går fra 0 mot negative verdier, går $F$ ganske fort mot 0. Dermed er nok ikke modellen gyldig for så mange år før 1960. Etter ca. $t=30$ flater grafen til $F$ ut. Mange små tettsteder i Norge har i nyere tid hatt liten befolkningsendring, og dermed er det ikke umulig at modellens gyldighetsområde strekker seg langt over $t=20$.
\item $f'(12)\approx 115$, som betyr at ved starten av 1972 økte antall innbyggere med ca. 115 personer per år. $f''(12)\approx-17$, som betyr at økningen i antall innbyggere var avtagende i starten av 1972.
\item Vi definerer $f'(x)$ som en funksjon, og finner skjæringspunktene med linja $y=150$. I ca. perioden 1963-1970 økte befolkningen med mer enn 150 personer per år.
}
\begin{figure}
	\centering
	\includegraphics[scale=0.2]{opg1}
\end{figure}

\newpage
\subsection*{Oppgave 2}
\abc{
\item Vi har at
\begin{align*}
	\lim\limits_{x\to-2^-} f(x)&=2\cdot(-2)-2 = -6 \vn
	\lim\limits_{x\to-2^+} f(x)&=2\cdot(-2)^3+2\cdot(-2)^2-2\cdot(-2)=-4
\end{align*}
Dermed er ikke $f$ kontinuerlig hvis $a=2$ og $b=-2$.
\item Vi setter $g(x)=ax+b$ og $h(x)=2x^3+2x^2-2x$. Krav at $f$ skal være kontinuerlig og deriverbar gir lilningssettene
\begin{align*}
	h'(k)=0 \\
	h(k)=c
\end{align*}
og 
\begin{align*}
	g(-2)=h(-2) \vn
	g'(-2)=h'(-2)
\end{align*}
}
I CAS finner vi at $a=15$ og $b=24$. I tillegg er enten $c=2$ og $k=-1$, eller så er $c=-\frac{10}{27}$ og $k=\frac{1}{3}$.
\begin{figure}
	\centering
	\includegraphics[scale=0.4]{opg2}
\end{figure}

\newpage
\subsection*{Oppgave 3}
\abc{
\item Vi definerer funksjonen i CAS, og finner at $l(500)\approx3,48$ og $l(1400)\approx4,1$. Altså svinger verdiene mellom det som tilsier at lukt bør begrenses og at tiltak kreves. Dermed har beboerne grunn til å klage.
\item For at lukt skal være akseptabel, må vi ha at $l=2$ eller mindre. Da kan ikke $C$ overstige ca. 44.
}
\begin{figure}
	\centering
	\includegraphics[scale=0.4]{opg3}
\end{figure}
\newpage
\subsection*{Oppgave 4}
\abc{
\item $\vv{HU}=[1200, 100]$. Siden dette er en retningsvektor for $l$, er
\[ l(s)=(0, 300)+s(1200, 100) \]
$l(0)=H$ og $l(1)=U$.
\item Siden Ina går med konstant fart, har hun etter $5\enh{min}$ gått $\frac{5}{20}$ av turen, og dermed $\frac{5}{20}$ av lengden til $\vv{HU}$. Da har hun gått ca. 301\enh{m} (celle 4).
\item I punktet hvor Jonas og Ina møtes er likningssettet fra celle 5 og cell 6 oppfylt. Ina møter Jonas når hun har gått $\frac{7}{20}$ av turen (celle 7), og da hun gått ca. 421\enh{m} (celle 8).
}
\begin{figure}
	\centering
	\includegraphics[scale=0.4]{opg4}
\end{figure}
\newpage
\subsection*{Oppgave 5}
\abc{
\item Vi har at 
\begin{align*}
	{\vec{p}\,}^2 &= \vec{a}^2+2\vec{a}\cdot\vec{b} +{\vec{b}\,}^2 \\
	&=4^2 +2\cdot4\cdot2\sqrt{3}\cdot \frac{\sqrt{3}}{2}+ (2\sqrt{3})^2 \\
	&=52
\end{align*}
Dermed er $|p|=\sqrt{52}$.
\item Vi har at
\begin{align*}
	\vec{p}\cdot\vec{q}&=0 \\
	(\vec{a}+\vec{b}) \cdot (t\vec{a}+ \vec{b})&=0 \\
	t\vec{a}^2+t\vec{a}\cdot\vec{b}+\vec{a}\cdot\vec{b}+{\vec{b}\,}^2 &=0 \\
	16t+12t+12+12&=0 \\
	t &= -\frac{6}{7}
\end{align*}
}

\newpage
\subsubsection{Oppgave 6}
\abc{
\item Vi skriver par med funksjoner først og andrederivert etter.\os
A og G. Siden $(a^x)'=a^x\ln x$, ligger grafen i G litt over grafen i A. \os
E og H. Av symmetrien viser grafen i E en andregradsfunksjon. Da er den dobbeltderiverte en konstant. (E kunne også vært andrederivert for D)
\os
B og C. Grafen i B går mot $\pm\infty$, og har et infleksjonspunkt. Da er dette en tredjegradsfunksjon, med en rett linje som dobbeltderivert. \os
F og D. Av symmetrien viser D en andregradsfunksjon. Da er den tilhørende funksjonen en fjerdegradsfunksjon. Ved å skrive $x^4$ i Geogebra ser vi at dette stemmer godt overens med F.
\item A, B, C og G. Dette fordi de er injektive.	
}

\newpage


\end{document}